% =====================================================================
% Figure contract
% 核心结论：十二个研究问题按研究逻辑线性递进，并直接映射到综述章节。
% 图原型：中等复杂度双行 pipeline；上行 RQ1--RQ6，下行 RQ7--RQ12。
% 证据层级：hero = 贯穿两行的 RQ chain；支撑 = 五类 pastel 分组与章节映射。
% 能不能更少？：不能；十二个 RQ、五个组名与全部章节映射均由规格明确要求。
%
% Step ① 设计文档（Form A / 中等复杂度结构图）
% 领域：Computer-Use Agents 中文综述；格式：TikZ standalone + XeLaTeX。
% W_print = 24 cm；W_ink 目标约 22.1 cm；目标印刷最小字号 7.0 pt。
% 布局：两行各六个 3.30 cm 节点，同行边距 0.45 cm，行间 corridor 1.45 cm。
% 连线：同行均为 3 pt 小 tip 短箭头；RQ6→RQ7 走行间正交 corridor。
% 视觉强调：五组低饱和浅底色；RQ 号粗体，章节映射灰色小字。
% 可视化决策：输入为纯结构映射，无数值，不嵌入 chart/heatmap。
%
% ASCII 草图：
%   [定义]       [能力]                                      [学习与基础设施]
%   RQ1 -> RQ2 -> RQ3 -> RQ4 -> RQ5 -> RQ6
%     +-------------------------------------------+
%     v
%   RQ7 -> RQ8 -> RQ9 -> RQ10 -> RQ11 -> RQ12
%   [学习与基础设施] [评测与可靠]             [部署]
%
% Pre-flight：同排节点由 positioning 构造；全部箭头端于 node anchor；
% 长横段位于两行之间；无中文旋转；直线不使用 rounded corners。
% =====================================================================

\documentclass[border=8pt]{standalone}
\usepackage{fontspec}
\usepackage{xeCJK}
\usepackage{xcolor}
\usepackage{tikz}
\setmainfont{Helvetica Neue}
\setsansfont{Helvetica Neue}
\setCJKmainfont{Heiti SC}
\setCJKsansfont{Heiti SC}
\usetikzlibrary{arrows.meta,bending,positioning,calc}

% Low-saturation palette shared with cua-survey-fig1-overview.tex.
\definecolor{ink}{HTML}{27364A}
\definecolor{muted}{HTML}{748092}
\definecolor{mainLine}{HTML}{4F667C}
\definecolor{blueLine}{HTML}{6D879E}
\definecolor{blueFill}{HTML}{E7EEF4}
\definecolor{purpleLine}{HTML}{817594}
\definecolor{purpleFill}{HTML}{EEEAF2}
\definecolor{goldLine}{HTML}{9B875B}
\definecolor{goldFill}{HTML}{F2EEE3}
\definecolor{tealLine}{HTML}{648D8A}
\definecolor{tealFill}{HTML}{E5F0EE}
\definecolor{roseLine}{HTML}{9A777B}
\definecolor{roseFill}{HTML}{F3E9EA}

\tikzset{
  every node/.style={
    font=\sffamily\fontsize{8.8}{10.5}\selectfont,
    text=ink
  },
  rqnode/.style={
    rectangle,
    rounded corners=4pt,
    line width=0.75pt,
    minimum width=3.30cm,
    minimum height=1.85cm,
    text width=2.92cm,
    align=center,
    inner xsep=5pt,
    inner ysep=5pt
  },
  group label/.style={
    rounded corners=2pt,
    text=white,
    inner xsep=6pt,
    inner ysep=2.4pt,
    font=\sffamily\fontsize{7.6}{9.0}\selectfont\bfseries
  },
  arrow short/.style={
    -{Stealth[length=3pt 1.0,width'=0pt 0.5,sep=0pt 0.3]},
    line width=0.8pt,
    shorten >=1pt,
    shorten <=0pt,
    color=mainLine
  },
  transition/.style={
    -{Stealth[length=5pt 1.5,width'=0pt 0.6,sep=0pt 0.5]},
    line width=0.9pt,
    shorten >=2pt,
    shorten <=1pt,
    color=mainLine
  }
}

\newcommand{\RQcontent}[3]{%
  {\fontsize{9.6}{11.0}\selectfont\bfseries #1}\\[2.5pt]
  {\fontsize{8.7}{10.2}\selectfont #2}\\[2.5pt]
  {\color{muted}\fontsize{7.4}{8.8}\selectfont #3}%
}

\begin{document}
\sffamily
\begin{tikzpicture}

% Upper row: RQ1--RQ6.
\node[rqnode,fill=blueFill,draw=blueLine] (rq1)
  {\RQcontent{RQ1}{定义与坐标}{§2–3}};
\node[rqnode,right=0.45cm of rq1,fill=purpleFill,draw=purpleLine] (rq2)
  {\RQcontent{RQ2}{感知}{§4.5·6.7}};
\node[rqnode,right=0.45cm of rq2,fill=purpleFill,draw=purpleLine] (rq3)
  {\RQcontent{RQ3}{动作接口}{§4.6–4.8}};
\node[rqnode,right=0.45cm of rq3,fill=purpleFill,draw=purpleLine] (rq4)
  {\RQcontent{RQ4}{模型架构}{§6.1–6.6}};
\node[rqnode,right=0.45cm of rq4,fill=purpleFill,draw=purpleLine] (rq5)
  {\RQcontent{RQ5}{规划与记忆}{§6.8–6.9}};
\node[rqnode,right=0.45cm of rq5,fill=goldFill,draw=goldLine] (rq6)
  {\RQcontent{RQ6}{训练与 RL}{§7}};

% Lower row: RQ7--RQ12, aligned to the upper row by construction.
\node[rqnode,below=1.45cm of rq1,fill=goldFill,draw=goldLine] (rq7)
  {\RQcontent{RQ7}{数据与经验}{§5}};
\node[rqnode,right=0.45cm of rq7,fill=goldFill,draw=goldLine] (rq8)
  {\RQcontent{RQ8}{环境与 runtime}{§4.2–4.4·9.6}};
\node[rqnode,right=0.45cm of rq8,fill=tealFill,draw=tealLine] (rq9)
  {\RQcontent{RQ9}{评测与验证}{§8}};
\node[rqnode,right=0.45cm of rq9,fill=tealFill,draw=tealLine] (rq10)
  {\RQcontent{RQ10}{可靠性与恢复}{§6.10·10.5}};
\node[rqnode,right=0.45cm of rq10,fill=tealFill,draw=tealLine] (rq11)
  {\RQcontent{RQ11}{安全与监督}{§6.11·8.9·10.8}};
\node[rqnode,right=0.45cm of rq11,fill=roseFill,draw=roseLine] (rq12)
  {\RQcontent{RQ12}{部署与开放问题}{§9–10}};

% Title and compact group labels.
\node[
  anchor=south,
  font=\sffamily\fontsize{12.0}{14.0}\selectfont\bfseries,
  text=ink
] at ([yshift=1.08cm]$(rq1.north)!0.5!(rq6.north)$)
  {十二个研究问题:线性递进与章节映射};

\node[group label,fill=blueLine,anchor=south]
  at ([yshift=0.30cm]rq1.north) {定义};
\node[group label,fill=purpleLine,anchor=south]
  at ([yshift=0.30cm]$(rq2.north)!0.5!(rq5.north)$) {能力};
\node[group label,fill=goldLine,anchor=north]
  at ([yshift=-0.30cm]$(rq7.south)!0.5!(rq8.south)$) {学习与基础设施};
\node[group label,fill=tealLine,anchor=north]
  at ([yshift=-0.30cm]$(rq9.south)!0.5!(rq11.south)$) {评测与可靠};
\node[group label,fill=roseLine,anchor=north]
  at ([yshift=-0.30cm]rq12.south) {部署};

% Short adjacent arrows: straight segments deliberately have no rounded corners.
\draw[arrow short] (rq1.east) -- (rq2.west);
\draw[arrow short] (rq2.east) -- (rq3.west);
\draw[arrow short] (rq3.east) -- (rq4.west);
\draw[arrow short] (rq4.east) -- (rq5.west);
\draw[arrow short] (rq5.east) -- (rq6.west);

\draw[arrow short] (rq7.east) -- (rq8.west);
\draw[arrow short] (rq8.east) -- (rq9.west);
\draw[arrow short] (rq9.east) -- (rq10.west);
\draw[arrow short] (rq10.east) -- (rq11.west);
\draw[arrow short] (rq11.east) -- (rq12.west);

% RQ6 -> RQ7: an orthogonal fold whose horizontal segment stays in the row gap.
\coordinate (transitionRight) at ([yshift=-0.72cm]rq6.south);
\coordinate (transitionLeft) at (rq7.north |- transitionRight);
\draw[transition,rounded corners=5pt]
  (rq6.south) -- (transitionRight) -- (transitionLeft) -- (rq7.north);

\end{tikzpicture}
\end{document}
